% Source: https://artofproblemsolving.com/community/c3h2434748_200th_post_problems
% Stable identity: aops-topic-2434748-q08
Let the squares $S_1$, $S_2$, $S_3$ have sides $s_1$, $s_2$, $s_3$.

\begin{asy}
usepackage("amsmath");
size(6.5cm);
defaultpen(fontsize(10pt));
fill((0,1)--(2,1)--(2,3)--(0,3)--cycle,rgb(0.85,0.92,1)); fill((2,1)--(3,1)--(3,2)--(2,2)--cycle,rgb(1,0.88,0.7)); fill((3,0)--(5,0)--(5,2)--(3,2)--cycle,rgb(0.85,1,0.85));
draw((0,0)--(5,0)--(5,3)--(0,3)--cycle,linewidth(0.9));
draw((3,0)--(3,1)--(0,1),linewidth(0.9));
draw((3,1)--(3,2)--(5,2),linewidth(0.9));
draw((3,2)--(2,2)--(2,1)--(2,3),linewidth(0.9));
label("$R_1$",(1.5,0.5)); label("$S_3$",(4,0.7)); label("$S_2$",(2.5,1.65)); label("$S_1$",(1,2.1)); label("$R_2$",(3.6,2.5));
draw((0,1.3)--(5,1.3),red+dashed);
label("$s_1$",(1,3),N); label("$s_2$",(2.5,2),N,fontsize(9pt)); label("$s_3$",(4,0),S); label("$s_3$",(5,1),E);
label("$s_1-s_2$",(5,2.5),E);
draw((0,-0.55)--(5,-0.55),Arrows(4)); label("$3322$",(2.5,-0.55),S);
draw((-0.6,0)--(-0.6,3),Arrows(4)); label("$2020$",(-0.6,1.5),W);
\end{asy}

\textbf{Across.} The squares $S_1$, $S_2$, $S_3$ stand side by side and span the whole width (the red dashed line crosses $S_1$, $S_2$, $S_3$ and nothing else), so
\[s_1+s_2+s_3=3322.\]
\textbf{Up the right side.} $S_1$ and $S_2$ both stand on top of $R_1$, so their bottom sides are at the same height; the top of $S_1$ is the top of the big rectangle, and $R_2$ lies between the top of $S_2$ and the top of the big rectangle. So $R_2$ has height $s_1-s_2$. Along the right side of the big rectangle we meet $S_3$ and then $R_2$, so
\[s_3+(s_1-s_2)=2020.\]
Subtracting the second equation from the first gives $2s_2=3322-2020=1302$, so $s_2=651$. The answer is $\boxed{\textbf{(A) }651}$.

(Such a tiling exists, for example with $s_1=s_3=\dfrac{2671}{2}$: then $R_1$ is $\dfrac{3973}{2}$ by $\dfrac{1369}{2}$ and $R_2$ is $\dfrac{3973}{2}$ by $\dfrac{1369}{2}$; the data determine $s_2$ but not $s_1$ and $s_3$ separately.)
